\documentclass[aspectratio=169,11pt]{beamer}
\usepackage[utf8]{inputenc}
\usepackage[T5]{fontenc}
\usepackage[vietnamese]{babel}
\usepackage{lmodern}
\usepackage{graphicx,tikz,amsmath}
\usetikzlibrary{arrows.meta,positioning,calc,fit,backgrounds}
\definecolor{BK}{RGB}{30,78,121}
\definecolor{Teal}{RGB}{15,120,123}
\definecolor{Amber}{RGB}{170,100,25}
\definecolor{Red}{RGB}{176,60,61}
\definecolor{Ink}{RGB}{34,48,62}
\definecolor{Muted}{RGB}{87,105,122}
\setbeamercolor{normal text}{fg=Ink,bg=white}
\setbeamercolor{structure}{fg=BK}
\setbeamerfont{frametitle}{size=\large,series=\bfseries}
\setbeamertemplate{navigation symbols}{}
\setbeamersize{text margin left=.65cm,text margin right=.65cm}
\setbeamertemplate{headline}{}
\setbeamertemplate{frametitle}{\vspace{.19cm}\strut\insertframetitle\par\vspace{.08cm}{\color{BK!22}\hrule height .6pt}\vspace{.06cm}}
\setbeamertemplate{footline}{\hspace{.65cm}\color{Muted}\fontsize{7}{8}\selectfont CLOUD COMPUTING\quad /\quad PHẦN 3: GPU SHARING \& OPTIMIZATION\hfill November 27\textsuperscript{th} 2026\qquad\insertframenumber\,/\,8\hspace{.65cm}\vspace{.17cm}}
\setbeamertemplate{background canvas}{\begin{tikzpicture}[remember picture,overlay]\node[opacity=.12,inner sep=0] at (current page.center){\includegraphics[width=\paperwidth,height=\paperheight]{assets/Background.pdf}};\fill[white] (current page.north west) rectangle ([yshift=-.92cm]current page.north east);\end{tikzpicture}}
\hypersetup{pdftitle={Cloud Computing - Phần 3: GPU Sharing và Resource Optimization},pdfauthor={Trần Nguyễn Việt Hưng; Nguyễn Hoàng Gia Huy; Nguyễn Đình Minh},colorlinks=true,urlcolor=BK}
\tikzset{
 box/.style={draw=BK!25,fill=white,rounded corners=3pt,line width=.7pt,inner sep=9pt,align=left,text=Ink,font=\fontsize{10}{12}\selectfont},
 smallbox/.style={box,font=\fontsize{9.5}{11.5}\selectfont,inner sep=7pt},
 arr/.style={-{Latex[length=2.2mm]},line width=1pt,draw=BK},
 softarr/.style={arr,draw=Muted!65},
 label/.style={font=\fontsize{10}{12}\selectfont,text=Muted,align=left},
 head/.style={font=\fontsize{12}{15}\selectfont\bfseries,text=BK,align=left},
 big/.style={font=\fontsize{27}{30}\selectfont\bfseries,text=BK},
 chip/.style={rounded corners=2pt,fill=BK!7,text=BK,inner sep=5pt,font=\fontsize{9}{11}\selectfont\bfseries},
 gpuchip/.style={rounded corners=1.5pt,fill=Teal!12,draw=Teal!45,text=Teal,inner xsep=4pt,inner ysep=2pt,font=\fontsize{7.5}{9}\selectfont\bfseries},
 scopebox/.style={draw=BK!30,dashed,rounded corners=4pt,line width=.75pt,fill=BK!2},
 }
\newenvironment{canvas}{\begin{tikzpicture}[x=1cm,y=1cm]\useasboundingbox (0,0) rectangle (14.7,6.35);}{\end{tikzpicture}}
% Varied bottom/callout helpers so slides do not repeat the same dark bar.
\newcommand{\accentbar}[3]{%
 \fill[#1!6,rounded corners=2.5pt] (0,.12) rectangle (14.55,.82);
 \fill[#1] (0,.12) rectangle (.13,.82);
 \node[anchor=west,text=Ink,font=\fontsize{10}{12}\selectfont,text width=13.95cm] at (.32,.47){\textcolor{#1}{\textbf{#2:}} #3};
}
\newcommand{\keynote}[2]{%
 \draw[BK!20,line width=.6pt] (0,.76)--(14.55,.76);
 \node[chip,fill=Teal!12,text=Teal,anchor=west,inner xsep=6pt,inner ysep=3pt] (knbadge) at (0,.38){#1};
 \path let \p1=(knbadge.east) in \pgfextra{\pgfmathsetlengthmacro{\keywidth}{14.55cm-\x1-.28cm}\xdef\KeyWidth{\keywidth}};
 \node[anchor=west,text=Ink,font=\fontsize{10}{12}\selectfont,text width=\KeyWidth] at ([xshift=.28cm]knbadge.east){#2};
}
% Audit actual node bounds, excluding enclosing/background shapes.
\newcommand{\auditnode}[2]{\path let \p1=(#2.south west),\p2=(#2.north east) in \pgfextra{\typeout{LAYOUT|\insertframenumber|#1|#2|\x1|\y1|\x2|\y2}};}
\newcommand{\closingbanner}[1]{%
 \node[box,fill=BK,draw=BK,text=white,text width=13.95cm,minimum height=.68cm,inner ysep=5pt,anchor=south west,font=\fontsize{10.2}{12}\selectfont\bfseries] at (0,.1){#1};
}
\providecommand{\SlideLimit}{8}
\begin{document}
% Slide 1 -- Part 3 title only, matching Cloud_Part1/slides/02-part1-title.tex.
\begin{frame}[plain]
\begin{tikzpicture}[remember picture,overlay]
\fill[BK] (current page.north west) rectangle ([xshift=.38cm]current page.south west);
\node[align=left,anchor=west,text=BK,font=\fontsize{30}{37}\selectfont\bfseries,text width=13.3cm] at ([xshift=1.3cm]current page.west){GPU Sharing,\\Resource Optimization\\\& Putting It Together};
\end{tikzpicture}
\end{frame}

\ifnum\SlideLimit>1 % Slide 2 -- Part A: GPU Underutilization & 4 Optimization Goals.
\begin{frame}[t]{Vì sao GPU bị lãng phí và bốn mục tiêu tối ưu}
\begin{canvas}
% Top-left: Single physical GPU underutilization visual
\node[box,fill=BK!3,text width=6.4cm,minimum height=2.34cm,anchor=north west] (leftcard) at (0,6.24){};
\node[head,anchor=north west] at (.24,6.10){Cấp nguyên 1 GPU cho workload nhỏ};
\node[label,text width=6.05cm,anchor=north west,font=\fontsize{8.6}{10.6}\selectfont] at (.24,5.58){Kubernetes chỉ nhận request GPU theo số nguyên (\texttt{nvidia.com/gpu: 1}).};

% Visual utilization bar inside left card
\fill[Teal!78,rounded corners=2pt] (.28,4.24) rectangle (2.15,4.82);
\node[anchor=center,text=white,font=\fontsize{8.2}{10}\selectfont\bfseries] at (1.215,4.53){Đang dùng};
\draw[draw=Red!75,fill=Red!6,dashed,line width=.9pt,rounded corners=2pt] (2.25,4.24) rectangle (6.65,4.82);
\node[anchor=center,text=Red,font=\fontsize{8.2}{10}\selectfont\bfseries] at (4.45,4.53){Compute \& VRAM nhàn rỗi};

% Top-right: 3 root causes of underutilization
\node[box,fill=Teal!4,draw=Teal!30,text width=6.4cm,minimum height=2.34cm,anchor=north west] (rightcard) at (7.6,6.24){};
\node[head,text=Teal,anchor=north west] at (7.84,6.10){Nguyên nhân chính gây Underutilization};
\node[anchor=north west,text=Ink,text width=6.15cm,font=\fontsize{8.7}{11.2}\selectfont] at (7.84,5.54){%
 $\bullet$\quad \textbf{Cấp phát nguyên chiếc:} workload nhỏ chiếm trọn 1 GPU.\\
 $\bullet$\quad \textbf{Tải không đều (bursty):} inference nhẹ, notebook có khoảng nghỉ.\\
 $\bullet$\quad \textbf{Nghẽn ngoài GPU:} chờ CPU tiền xử lý hoặc I/O đọc dữ liệu.};

% Middle heading + 4 optimization goals (matching Slide 4 of Part 1)
\node[head,anchor=north west,inner sep=0] (goalheading) at (0,3.68){Bốn mục tiêu khi tối ưu tài nguyên GPU};

\foreach \x/\col/\ttl/\body [count=\i] in {
 0/Teal/Utilization/{Tăng tỷ lệ sử dụng thực tế của compute và VRAM.},
 3.78/BK/Performance/{Giữ throughput cao và kiểm soát độ trễ (latency).},
 7.56/Amber/Isolation/{Cô lập bộ nhớ, hiệu năng và lỗi giữa các workload.},
 11.34/Red/Cost/{Giảm số GPU cần cấp, tối ưu chi phí vận hành.}
}{
 \node[smallbox,text width=2.82cm,minimum height=1.96cm,anchor=north west,inner sep=7pt] (goal\i) at (\x,2.96){};
 \fill[\col] (\x,2.96) rectangle (\x+3.32,2.88);
 \node[anchor=north west,text=\col,text width=3cm,font=\fontsize{9.5}{11.5}\selectfont\bfseries] at (\x+.14,2.72){\ttl};
 \node[anchor=north west,text=Ink,text width=2.99cm,font=\fontsize{8.4}{10.5}\selectfont] at (\x+.14,2.18){\body};
}
\auditnode{cards}{leftcard}
\auditnode{cards}{rightcard}
\auditnode{goals}{goalheading}
\auditnode{goals}{goal1}
\auditnode{goals}{goal2}
\auditnode{goals}{goal3}
\auditnode{goals}{goal4}
\accentbar{Teal}{Bài toán cốt lõi}{Tối ưu GPU luôn là sự đánh đổi giữa Utilization, Performance, Isolation và Cost.}
\end{canvas}
\end{frame}
\fi
\ifnum\SlideLimit>2 % Slide 3 -- Part B1: GPU Time-Slicing (Time-Sharing).
\begin{frame}[t]{GPU Time-Slicing: Chia sẻ GPU theo lát cắt thời gian}
\begin{canvas}
\node[chip,anchor=north west,inner ysep=2.5pt,font=\fontsize{8.2}{10}\selectfont\bfseries] at (0,6.28){KỸ THUẬT 1\quad /\quad TIME-SLICING (GKE TIME-SHARING)};

% Top diagram: 3 Pods -> Sequential Time-Sliced GPU Timeline
\node[scopebox,draw=BK!35,fill=BK!2,minimum width=14.55cm,minimum height=2.04cm,anchor=north west] (tsframe) at (0,5.78){};

\node[smallbox,fill=BK!8,draw=BK!45,text width=1.65cm,align=center,minimum height=.44cm,inner sep=3pt,font=\fontsize{8.5}{10}\selectfont\bfseries] (pa) at (1.15,5.38){Pod A};
\node[smallbox,fill=Teal!10,draw=Teal!50,text width=1.65cm,align=center,minimum height=.44cm,inner sep=3pt,font=\fontsize{8.5}{10}\selectfont\bfseries] (pb) at (1.15,4.76){Pod B};
\node[smallbox,fill=Amber!10,draw=Amber!50,text width=1.65cm,align=center,minimum height=.44cm,inner sep=3pt,font=\fontsize{8.5}{10}\selectfont\bfseries] (pc) at (1.15,4.14){Pod C};

\node[smallbox,fill=white,draw=BK!30,text width=10.35cm,minimum height=1.66cm,anchor=north west,inner sep=6pt] (timeline) at (3.45,5.59){};
\node[anchor=north west,text=BK,font=\fontsize{9}{10.5}\selectfont\bfseries] at (3.65,5.47){1 GPU vật lý --- Chuyển ngữ cảnh luân phiên (Context Switching)};

% Alternating time slices
\foreach \x/\col/\lbl in {
 3.68/BK/A, 5.40/Teal/B, 7.12/Amber/C, 8.84/BK/A, 10.56/Teal/B, 12.28/Amber/C
}{
 \fill[\col!75,rounded corners=2pt] (\x,4.52) rectangle (\x+1.52,5.02);
 \node[anchor=center,text=white,font=\fontsize{8.5}{10}\selectfont\bfseries] at (\x+.76,4.77){Pod \lbl};
}
\draw[softarr] (3.68,4.28)--(13.85,4.28);
\node[anchor=west,text=Muted,font=\fontsize{8}{9.5}\selectfont] at (3.68,4.11){Thời gian (Instruction-level Preemption từ kiến trúc Pascal+)};

\draw[arr] (pa.east)--(timeline.west |- pa.east);
\draw[arr,draw=Teal] (pb.east)--(timeline.west |- pb.east);
\draw[arr,draw=Amber] (pc.east)--(timeline.west |- pc.east);

% Middle 2-column structured cards
\node[smallbox,fill=BK!3,draw=BK!30,text width=6.5cm,minimum height=2.40cm,anchor=north west] (tspro) at (0,3.48){};
\node[anchor=north west,text=BK,font=\fontsize{10.2}{12}\selectfont\bfseries] at (.22,3.30){Nguyên lý \& Đặc điểm trên GKE};
\node[anchor=north west,text=Ink,text width=6.25cm,font=\fontsize{8.7}{11.4}\selectfont] at (.22,2.78){%
 $\bullet$\quad Cấp luân phiên lát cắt thời gian trên cùng 1 GPU vật lý.\\
 $\bullet$\quad Mỗi container thấy toàn bộ GPU khi tới lượt thực thi.\\
 $\bullet$\quad Hỗ trợ rộng (Pascal+); GKE cho tối đa \textbf{48 client/GPU}.};

\node[smallbox,fill=Amber!5,draw=Amber!40,text width=6.5cm,minimum height=2.40cm,anchor=north west] (tscon) at (7.55,3.48){};
\node[anchor=north west,text=Amber,font=\fontsize{10.2}{12}\selectfont\bfseries] at (7.77,3.30){Đánh đổi \& Giới hạn cô lập};
\node[anchor=north west,text=Ink,text width=6.25cm,font=\fontsize{8.7}{11.4}\selectfont] at (7.77,2.78){%
 $\bullet$\quad Chi phí context switching làm tăng và dao động độ trễ.\\
 $\bullet$\quad \textbf{Không ép giới hạn VRAM} giữa các job (dễ gây OOM).\\
 $\bullet$\quad Chỉ nên dùng chung trong cùng \textit{trust boundary}.};

\auditnode{cards}{timeline}
\auditnode{cards}{tspro}
\auditnode{cards}{tscon}
\keynote{Phù hợp}{Workload dạng bursty, có khoảng nghỉ (dev/test, notebook, thử nghiệm) chấp nhận biến động latency.}
\end{canvas}
\end{frame}
\fi
\ifnum\SlideLimit>3 % Slide 4 -- Part B2: NVIDIA MPS (Multi-Process Service).
\begin{frame}[t]{NVIDIA MPS: Thực thi song song nhiều tiến trình CUDA}
\begin{canvas}
\node[chip,anchor=north west,inner ysep=2.5pt,font=\fontsize{8.2}{10}\selectfont\bfseries] at (0,6.28){KỸ THUẬT 2\quad /\quad NVIDIA MULTI-PROCESS SERVICE (MPS)};

% Top diagram: Clients -> MPS Server -> Concurrent GPU Execution
\node[scopebox,draw=Teal!45,fill=Teal!2,minimum width=14.55cm,minimum height=2.04cm,anchor=north west] (mpsframe) at (0,5.78){};

\node[smallbox,fill=BK!8,draw=BK!45,text width=1.85cm,align=center,minimum height=.44cm,inner sep=3pt,font=\fontsize{8.4}{10}\selectfont\bfseries] (ca) at (1.25,5.38){CUDA Client A};
\node[smallbox,fill=Teal!10,draw=Teal!50,text width=1.85cm,align=center,minimum height=.44cm,inner sep=3pt,font=\fontsize{8.4}{10}\selectfont\bfseries] (cb) at (1.25,4.76){CUDA Client B};
\node[smallbox,fill=Amber!10,draw=Amber!50,text width=1.85cm,align=center,minimum height=.44cm,inner sep=3pt,font=\fontsize{8.4}{10}\selectfont\bfseries] (cc) at (1.25,4.14){CUDA Client C};

\node[smallbox,fill=white,draw=Teal!50,text width=3.1cm,minimum height=1.56cm,align=center,inner sep=6pt] (mpssrv) at (5.35,4.76){%
 \textcolor{Teal}{\textbf{MPS Server}}\\[.08cm]
 {\fontsize{8}{9.8}\selectfont Triển khai CUDA API\\tương thích nhị phân\\Gộp luồng lập lịch}};

\node[smallbox,fill=white,draw=BK!35,text width=5.25cm,minimum height=1.66cm,anchor=north west,inner sep=6pt] (mpsgpu) at (8.55,5.59){};
\node[anchor=north west,text=BK,font=\fontsize{8.8}{10.5}\selectfont\bfseries] at (8.74,5.47){1 GPU vật lý --- Chạy song song};
\fill[BK!75,rounded corners=2pt] (8.76,4.90) rectangle (13.98,5.16);
\node[anchor=center,text=white,font=\fontsize{7.8}{9}\selectfont\bfseries] at (11.37,5.03){Kernels Client A (SMs)};
\fill[Teal!75,rounded corners=2pt] (8.76,4.54) rectangle (13.98,4.80);
\node[anchor=center,text=white,font=\fontsize{7.8}{9}\selectfont\bfseries] at (11.37,4.67){Kernels Client B (SMs)};
\fill[Amber!80,rounded corners=2pt] (8.76,4.18) rectangle (13.98,4.44);
\node[anchor=center,text=white,font=\fontsize{7.8}{9}\selectfont\bfseries] at (11.37,4.31){Kernels Client C (SMs)};

\draw[arr] (ca.east)--(mpssrv.west |- ca.east);
\draw[arr,draw=Teal] (cb.east)--(mpssrv.west);
\draw[arr,draw=Amber] (cc.east)--(mpssrv.west |- cc.east);
\draw[arr,draw=Teal] (mpssrv.east)--(mpsgpu.west |- mpssrv.east);

% Middle 2-column structured cards
\node[smallbox,fill=Teal!4,draw=Teal!35,text width=6.5cm,minimum height=2.40cm,anchor=north west] (mpspro) at (0,3.48){};
\node[anchor=north west,text=Teal,font=\fontsize{10.2}{12}\selectfont\bfseries] at (.22,3.30){Nguyên lý \& Điểm mạnh};
\node[anchor=north west,text=Ink,text width=6.25cm,font=\fontsize{8.7}{11.4}\selectfont] at (.22,2.78){%
 $\bullet$\quad Các CUDA kernel chạy \textbf{đồng thời} để lấp đầy SM của GPU.\\
 $\bullet$\quad Giảm overhead context switching và bộ nhớ ngữ cảnh.\\
 $\bullet$\quad GKE giới hạn mềm \textbf{\% active threads} \& \textbf{pinned memory}.};

\node[smallbox,fill=Red!5,draw=Red!35,text width=6.5cm,minimum height=2.40cm,anchor=north west] (mpscon) at (7.55,3.48){};
\node[anchor=north west,text=Red,font=\fontsize{10.2}{12}\selectfont\bfseries] at (7.77,3.30){Hạn chế \& Rủi ro cô lập};
\node[anchor=north west,text=Ink,text width=6.25cm,font=\fontsize{8.7}{11.4}\selectfont] at (7.77,2.78){%
 $\bullet$\quad Không cô lập phần cứng cho memory bandwidth/encoder.\\
 $\bullet$\quad \textbf{Lỗi nghiêm trọng ở 1 client} lan sang mọi client trên GPU.\\
 $\bullet$\quad GKE yêu cầu \texttt{hostIPC: true} (cần chú ý bảo mật).};

\auditnode{cards}{mpssrv}
\auditnode{cards}{mpsgpu}
\auditnode{cards}{mpspro}
\auditnode{cards}{mpscon}
\keynote{Phù hợp}{Batch nhỏ/vừa hoặc các tiến trình hợp tác (cooperative MPI/CUDA) cần tối đa hóa throughput.}
\end{canvas}
\end{frame}
\fi
\ifnum\SlideLimit>4 % Slide 5 -- Part B3: Multi-Instance GPU (MIG).
\begin{frame}[t]{Multi-Instance GPU (MIG): Phân chia phần cứng độc lập}
\begin{canvas}
\node[chip,anchor=north west,inner ysep=2.5pt,font=\fontsize{8.2}{10}\selectfont\bfseries] at (0,6.28){KỸ THUẬT 3\quad /\quad MULTI-INSTANCE GPU (MIG)};

% Top diagram: Hardware-isolated GPU partitions (Ampere+)
\node[scopebox,draw=BK!40,fill=BK!2,minimum width=14.55cm,minimum height=2.04cm,anchor=north west] (migframe) at (0,5.78){};
\node[anchor=north west,text=BK,font=\fontsize{8.8}{10.5}\selectfont\bfseries] at (.22,5.67){1 GPU vật lý (Ampere+: A100, H100) --- Chia tối đa 7 Hardware Instances độc lập};

\node[smallbox,fill=Teal!7,draw=Teal!55,text width=3.95cm,minimum height=1.22cm,anchor=north west,inner sep=6pt] (m1) at (.24,5.24){%
 \textcolor{Teal}{\textbf{Instance 1 (1g.5gb)}}\\[.04cm]
 {\fontsize{8.1}{9.8}\selectfont SM + VRAM + BW riêng\\Gán trực tiếp cho \textbf{Pod A}}};

\node[smallbox,fill=BK!6,draw=BK!45,text width=4.20cm,minimum height=1.22cm,anchor=north west,inner sep=6pt] (m2) at (4.86,5.24){%
 \textcolor{BK}{\textbf{Instance 2 (2g.10gb)}}\\[.04cm]
 {\fontsize{8.1}{9.8}\selectfont SM + VRAM + BW riêng\\Gán trực tiếp cho \textbf{Pod B}}};

\node[smallbox,fill=Amber!7,draw=Amber!50,text width=4.25cm,minimum height=1.22cm,anchor=north west,inner sep=6pt] (m3) at (9.74,5.24){%
 \textcolor{Amber}{\textbf{Instance 3 (3g.20gb)}}\\[.04cm]
 {\fontsize{8.1}{9.8}\selectfont SM + VRAM + BW riêng\\Gán trực tiếp cho \textbf{Pod C}}};

% Hardware isolation wall indicators
\draw[draw=Ink!45,line width=1.1pt,dashed] (4.70,5.28)--(4.70,3.96);
\draw[draw=Ink!45,line width=1.1pt,dashed] (9.58,5.28)--(9.58,3.96);

% Middle 2-column structured cards
\node[smallbox,fill=Teal!4,draw=Teal!35,text width=6.5cm,minimum height=2.40cm,anchor=north west] (migpro) at (0,3.48){};
\node[anchor=north west,text=Teal,font=\fontsize{10.2}{12}\selectfont\bfseries] at (.22,3.30){Ưu điểm \& Cô lập phần cứng};
\node[anchor=north west,text=Ink,text width=6.25cm,font=\fontsize{8.7}{11.4}\selectfont] at (.22,2.78){%
 $\bullet$\quad \textbf{Hardware isolation:} tách biệt compute, VRAM, băng thông \& lỗi.\\
 $\bullet$\quad \textbf{QoS ổn định:} throughput và latency dự đoán được.\\
 $\bullet$\quad Có thể kết hợp MIG + Time-Slicing/MPS trên từng slice.};

\node[smallbox,fill=Amber!5,draw=Amber!40,text width=6.5cm,minimum height=2.40cm,anchor=north west] (migcon) at (7.55,3.48){};
\node[anchor=north west,text=Amber,font=\fontsize{10.2}{12}\selectfont\bfseries] at (7.77,3.30){Giới hạn \& Lưu ý trên GKE};
\node[anchor=north west,text=Ink,text width=6.25cm,font=\fontsize{8.7}{11.4}\selectfont] at (7.77,2.78){%
 $\bullet$\quad Chỉ hỗ trợ GPU dòng Ampere trở lên (A100, H100, ...).\\
 $\bullet$\quad Tối đa \textbf{7 instances/GPU}; phân vùng cố định kém linh hoạt.\\
 $\bullet$\quad GKE chọn qua \texttt{gke-gpu-partition-size}; chưa có \texttt{duty\_cycle}.};

\auditnode{cards}{m1}
\auditnode{cards}{m2}
\auditnode{cards}{m3}
\auditnode{cards}{migpro}
\auditnode{cards}{migcon}
\keynote{Phù hợp}{AI Inference production hoặc workload song song cần QoS ổn định và cô lập lỗi nghiêm ngặt.}
\end{canvas}
\end{frame}
\fi
\ifnum\SlideLimit>5 % Slide 6 -- Part B4: Comparison of Time-Slicing, NVIDIA MPS, and MIG.
\begin{frame}[t]{So sánh Time-Slicing, NVIDIA MPS và MIG}
\begin{canvas}
% Header row (y = 5.66 .. 6.24)
\fill[BK,rounded corners=2.5pt] (0,5.66) rectangle (2.68,6.24);
\node[anchor=west,text=white,font=\fontsize{9.2}{11}\selectfont\bfseries] at (.16,5.95){Tiêu chí};

\fill[Amber!90,rounded corners=2.5pt] (2.78,5.66) rectangle (6.62,6.24);
\node[anchor=west,text=white,font=\fontsize{9.2}{11}\selectfont\bfseries] at (2.94,5.95){Time-Slicing};

\fill[BK!82,rounded corners=2.5pt] (6.72,5.66) rectangle (10.56,6.24);
\node[anchor=west,text=white,font=\fontsize{9.2}{11}\selectfont\bfseries] at (6.88,5.95){NVIDIA MPS};

\fill[Teal,rounded corners=2.5pt] (10.66,5.66) rectangle (14.55,6.24);
\node[anchor=west,text=white,font=\fontsize{9.2}{11}\selectfont\bfseries] at (10.82,5.95){MIG (Multi-Instance)};

% 5 comparison rows with alternating subtle backgrounds
\foreach \y/\bg/\crit/\cOne/\cTwo/\cThree [count=\r] in {
 4.74/BK!3/Cơ chế chia sẻ/{Luân phiên lát cắt thời gian\\(Context Switching)}/{Chạy song song nhiều tiến\\trình CUDA qua MPS Server}/{Chia phần cứng thành\\tối đa 7 instances độc lập},
 3.82/white/Mức độ cô lập/{Phần mềm; \textbf{không ép VRAM}\\(nguy cơ OOM chéo)}/{Phần mềm (thread \%, VRAM);\\lỗi nghiêm trọng dễ lan}/{Phần cứng hoàn toàn\\(SM, VRAM, băng thông, lỗi)},
 2.90/BK!3/Hiệu năng \& QoS/{Có overhead chuyển ngữ\\cảnh; latency biến động}/{Throughput cao, giảm mạnh\\context switching}/{QoS ổn định, dự đoán được\\không nhiễu giữa các slice},
 1.98/white/GPU \& GKE/{Rộng (Pascal+); tối đa\\48 client/GPU vật lý}/{Volta+ tối ưu ($\le$ 48 client);\\cần \texttt{hostIPC: true}}/{Chỉ Ampere+ (A100, H100);\\tối đa 7 slice/GPU},
 1.06/BK!3/Workload phù hợp/{Dev/test, notebook, tải\\bursty có khoảng nghỉ}/{Batch nhỏ/vừa, tiến trình\\hợp tác (MPI/CUDA)}/{Inference production cần\\SLA và cô lập lỗi chặt}
}{
 \fill[\bg,rounded corners=2pt] (0,\y) rectangle (2.68,\y+.82);
 \draw[BK!18,rounded corners=2pt] (0,\y) rectangle (2.68,\y+.82);
 \node[anchor=west,text=BK,font=\fontsize{8.6}{10.2}\selectfont\bfseries,text width=2.42cm] at (.14,\y+.41){\crit};

 \fill[\bg,rounded corners=2pt] (2.78,\y) rectangle (6.62,\y+.82);
 \draw[BK!18,rounded corners=2pt] (2.78,\y) rectangle (6.62,\y+.82);
 \node[anchor=west,text=Ink,font=\fontsize{8.2}{10}\selectfont,text width=3.58cm] at (2.92,\y+.41){\cOne};

 \fill[\bg,rounded corners=2pt] (6.72,\y) rectangle (10.56,\y+.82);
 \draw[BK!18,rounded corners=2pt] (6.72,\y) rectangle (10.56,\y+.82);
 \node[anchor=west,text=Ink,font=\fontsize{8.2}{10}\selectfont,text width=3.58cm] at (6.86,\y+.41){\cTwo};

 \fill[\bg,rounded corners=2pt] (10.66,\y) rectangle (14.55,\y+.82);
 \draw[BK!18,rounded corners=2pt] (10.66,\y) rectangle (14.55,\y+.82);
 \node[anchor=west,text=Ink,font=\fontsize{8.2}{10}\selectfont,text width=3.62cm] at (10.80,\y+.41){\cThree};
}

\accentbar{Teal}{Kết hợp thực tế}{Trên GKE có thể chạy Time-Slicing hoặc MPS trên từng MIG partition để tăng mật độ container.}
\end{canvas}
\end{frame}
\fi
\ifnum\SlideLimit>6 % Slide 7 -- Part C1: Putting Everything Together (WHEN - WHERE - HOW).
\begin{frame}[t]{Tổng hợp kiến trúc quản lý và lập lịch GPU trên GKE}
\begin{canvas}
\node[label,text=Ink,anchor=north west,text width=14.5cm,font=\fontsize{9.6}{11.5}\selectfont] at (0,6.26){Ba lớp điều phối bổ trợ nhau từ lúc nộp job đến khi thực thi trên từng GPU vật lý.};

% Pillar 1: WHEN (Part 1)
\node[smallbox,fill=BK!3,draw=BK!35,text width=3.95cm,minimum height=4.42cm,anchor=north west] (pwhen) at (0,5.66){};
\fill[BK] (0,5.66) rectangle (4.45,5.56);
\node[chip,fill=BK,text=white,anchor=north west,inner xsep=7pt,inner ysep=3pt] at (.20,5.38){WHEN --- PHẦN 1};
\node[anchor=north west,text=BK,text width=4.05cm,font=\fontsize{10.5}{12.5}\selectfont\bfseries] at (.20,4.78){Kueue \& Gang Scheduling};
\node[anchor=north west,text=Ink,text width=4.05cm,font=\fontsize{8.6}{11.2}\selectfont] at (.20,3.96){%
 $\bullet$\quad Điều phối \textbf{thời điểm chạy} ở cấp Workload.\\[.12cm]
 $\bullet$\quad Hàng đợi, Quota, Priority, Preemption \& Fair Sharing.\\[.12cm]
 $\bullet$\quad \textbf{All-or-Nothing:} tránh giữ GPU nhàn rỗi khi thiếu worker.};

% Pillar 2: WHERE (Part 2)
\node[smallbox,fill=Teal!4,draw=Teal!40,text width=3.95cm,minimum height=4.42cm,anchor=north west] (pwhere) at (5.05,5.66){};
\fill[Teal] (5.05,5.66) rectangle (9.50,5.56);
\node[chip,fill=Teal,text=white,anchor=north west,inner xsep=7pt,inner ysep=3pt] at (5.25,5.38){WHERE --- PHẦN 2};
\node[anchor=north west,text=Teal,text width=4.05cm,font=\fontsize{10.5}{12.5}\selectfont\bfseries] at (5.25,4.78){Topology-Aware Scheduling};
\node[anchor=north west,text=Ink,text width=4.05cm,font=\fontsize{8.6}{11.2}\selectfont] at (5.25,3.96){%
 $\bullet$\quad Chọn \textbf{vị trí đặt} GPU/worker theo cấu trúc Node, Rack, Cluster.\\[.12cm]
 $\bullet$\quad Ưu tiên NVLink/NVSwitch và băng thông nội bộ cao.\\[.12cm]
 $\bullet$\quad Giảm độ trễ đồng bộ \textbf{AllReduce} khi training phân tán.};

% Pillar 3: HOW (Part 3)
\node[smallbox,fill=Amber!5,draw=Amber!45,text width=3.95cm,minimum height=4.42cm,anchor=north west] (phow) at (10.10,5.66){};
\fill[Amber] (10.10,5.66) rectangle (14.55,5.56);
\node[chip,fill=Amber,text=white,anchor=north west,inner xsep=7pt,inner ysep=3pt] at (10.30,5.38){HOW --- PHẦN 3};
\node[anchor=north west,text=Amber,text width=4.05cm,font=\fontsize{10.5}{12.5}\selectfont\bfseries] at (10.30,4.78){GPU Sharing \& Optimization};
\node[anchor=north west,text=Ink,text width=4.05cm,font=\fontsize{8.6}{11.2}\selectfont] at (10.30,3.96){%
 $\bullet$\quad Khai thác \textbf{hiệu quả từng GPU} cho workload nhỏ/inference.\\[.12cm]
 $\bullet$\quad Chọn Time-Slicing, NVIDIA MPS hoặc MIG theo ngữ cảnh.\\[.12cm]
 $\bullet$\quad Cân bằng Utilization, Performance, Isolation \& Cost.};

% Connecting arrows between pillars
\draw[arr,line width=1.15pt] (pwhen.east)--(pwhere.west);
\draw[arr,draw=Teal,line width=1.15pt] (pwhere.east)--(phow.west);

\auditnode{cards}{pwhen}
\auditnode{cards}{pwhere}
\auditnode{cards}{phow}
\closingbanner{Phối hợp WHEN -- WHERE -- HOW giúp tối ưu cả tiến độ huấn luyện lẫn chi phí hạ tầng GPU.}
\end{canvas}
\end{frame}
\fi
\ifnum\SlideLimit>7 % Slide 8 -- Part C2: Integrated Demo Introduction.
\begin{frame}[t]{Demo tổng hợp: AI Workload Scheduling với Kueue}
\begin{canvas}
\node[chip,anchor=north west,inner ysep=2.5pt,font=\fontsize{8.2}{10}\selectfont\bfseries] at (0,6.28){DEMO TỔNG HỢP\quad /\quad KUBERNETES + KUEUE};

% Objective callout box
\node[smallbox,fill=BK!4,draw=BK!30,text width=14.05cm,minimum height=1.12cm,anchor=north west,inner sep=7pt] (demogoal) at (0,5.76){};
\node[anchor=north west,text=BK,font=\fontsize{10}{12}\selectfont\bfseries] at (.22,5.60){Mục tiêu Demo: AI Workload Scheduling with Kueue on Kubernetes};
\node[anchor=north west,text=Ink,text width=13.95cm,font=\fontsize{8.9}{10.8}\selectfont] at (.22,5.16){Minh họa cách Kueue quản lý nhiều AI workloads đồng thời, kiểm soát quota theo hàng đợi và phối hợp đặt Pod theo topology trên cụm Kubernetes.};

% 3 Demo stages
\node[smallbox,fill=white,draw=BK!35,text width=3.95cm,minimum height=3.02cm,anchor=north west] (d1) at (0,4.30){};
\fill[BK] (0,4.30) rectangle (4.45,4.22);
\node[anchor=north west,text=BK,font=\fontsize{10}{12}\selectfont\bfseries] at (.20,4.04){1. Thiết lập hàng đợi};
\node[anchor=north west,text=Ink,text width=4.05cm,font=\fontsize{8.6}{11.2}\selectfont] at (.20,3.50){%
 $\bullet$\quad Khởi tạo \texttt{ResourceFlavor} và định nghĩa Topology.\\[.1cm]
 $\bullet$\quad Tạo \texttt{ClusterQueue} với hạn mức GPU cụ thể.\\[.1cm]
 $\bullet$\quad Gắn \texttt{LocalQueue} cho namespace chạy job.};

\node[smallbox,fill=white,draw=Teal!40,text width=3.95cm,minimum height=3.02cm,anchor=north west] (d2) at (5.05,4.30){};
\fill[Teal] (5.05,4.30) rectangle (9.50,4.22);
\node[anchor=north west,text=Teal,font=\fontsize{10}{12}\selectfont\bfseries] at (5.25,4.04){2. Nộp nhiều AI Workloads};
\node[anchor=north west,text=Ink,text width=4.05cm,font=\fontsize{8.6}{11.2}\selectfont] at (5.25,3.50){%
 $\bullet$\quad Nộp đồng thời các job có nhu cầu GPU khác nhau.\\[.1cm]
 $\bullet$\quad Quan sát quá trình \texttt{Queue} $\rightarrow$ \texttt{Admission} theo quota.\\[.1cm]
 $\bullet$\quad Kiểm chứng job vượt quota được giữ lại an toàn.};

\node[smallbox,fill=white,draw=Amber!45,text width=3.95cm,minimum height=3.02cm,anchor=north west] (d3) at (10.10,4.30){};
\fill[Amber] (10.10,4.30) rectangle (14.55,4.22);
\node[anchor=north west,text=Amber,font=\fontsize{10}{12}\selectfont\bfseries] at (10.30,4.04){3. Placement \& Kiểm chứng};
\node[anchor=north west,text=Ink,text width=4.05cm,font=\fontsize{8.6}{11.2}\selectfont] at (10.30,3.50){%
 $\bullet$\quad Quan sát vị trí Pod được gán theo miền topology.\\[.1cm]
 $\bullet$\quad Kiểm tra tài nguyên giải phóng khi job hoàn tất.\\[.1cm]
 $\bullet$\quad Job kế tiếp trong hàng đợi tự động được admit.};

\draw[arr,line width=1.1pt] (d1.east)--(d2.west);
\draw[arr,draw=Teal,line width=1.1pt] (d2.east)--(d3.west);

\auditnode{cards}{demogoal}
\auditnode{cards}{d1}
\auditnode{cards}{d2}
\auditnode{cards}{d3}
\keynote{Chuyển tiếp}{Mời thầy và các bạn theo dõi phần Demo trực tiếp trên cụm Kubernetes.}
\end{canvas}
\end{frame}
\fi
\end{document}
